\documentclass[11pt]{article}
\usepackage[margin=1in]{geometry}
\usepackage{amsmath,amssymb,amsthm,mathtools}
\usepackage{graphicx}
\usepackage[hidelinks]{hyperref}
\usepackage{xcolor}

\newtheorem{theorem}{Theorem}
\newtheorem{corollary}[theorem]{Corollary}
\newtheorem{lemma}[theorem]{Lemma}
\newcommand{\Sone}{\mathbb S^1}
\newcommand{\Ead}{\mathcal E_{a,d}}
\newcommand{\dist}{\operatorname{dist}}
\newcommand{\dd}{\,\mathrm d}

\title{Single-vortex extremizers on analytic planar domains\\
\large A partial result for Open Problem 1 of arXiv:1606.01526}
\author{Automated mathematics run \texttt{fa\_banach\_001}, lane 8}
\date{22 July 2026}

\begin{document}
\maketitle

\noindent\textbf{Status.}
Candidate substantial partial solution, likely valid, requiring human review.
The result below completely classifies attainment in every equivalence class
with one isolated vortex on a real-analytic simply connected planar domain.
It does not classify classes with several or non-atomic singularities.

\section{Source question}

The source is H.~Brezis, P.~Mironescu, and I.~Shafrir,
\emph{Distances between classes in $W^{1,1}(\Omega;\mathbb S^1)$},
arXiv:1606.01526, later published in \emph{Calc. Var. PDE} 57 (2018),
Article 14 \cite{BMS}. On page 22, Open Problem 1 asks whether some smooth
bounded domain fails property $(\mathrm P_1)$, namely the existence of
$u\in W^{1,1}(\Omega;\Sone)$ such that
\[
  0<\Sigma(u)=\int_\Omega |\nabla u|.
\]
The authors specifically suggest that the interior of a noncircular ellipse
may fail $(\mathrm P_1)$.

\begin{center}
\includegraphics[width=0.98\textwidth]{figures/open_problem_crop.png}
\end{center}

\section{Definitions and statement}

Let $\Omega\subset\mathbb R^2$ be a bounded simply connected domain, let
$a\in\Omega$, and let $d\in\mathbb Z\setminus\{0\}$. Define
\[
 w_{a,d}(x)=\left(\frac{x-a}{|x-a|}\right)^d,
 \qquad x\in\Omega\setminus\{a\}.
\]
Since $|\nabla w_{a,d}|=|d|/|x-a|$, this map belongs to
$W^{1,1}(\Omega;\Sone)$. Its equivalence class in the terminology of
\cite{BMS} is
\[
 \Ead=\{w_{a,d}e^{i\varphi}:\varphi\in W^{1,1}(\Omega;\mathbb R)\}.
\]
Every member of this class has one topological singularity at $a$, of
multiplicity $d$. Put
\[
 r_a=\dist(a,\partial\Omega),\qquad
 m_{a,d}=\inf_{u\in\Ead}\int_\Omega|\nabla u|.
\]
The geometric formula $\Sigma=2\pi L$ in \cite{BMS} gives
\begin{equation}\label{eq:source-value}
 m_{a,d}=\Sigma(w_{a,d})=2\pi |d|r_a.
\end{equation}

\begin{theorem}[single-vortex attainment classification]\label{thm:main}
Assume in addition that $\partial\Omega$ is connected and real analytic.
Then the infimum $m_{a,d}$ is attained by some $u\in\Ead$ if and only if
$\Omega=B_{r_a}(a)$.
\end{theorem}

\begin{corollary}[noncircular ellipses]\label{cor:ellipse}
Let $\Omega$ be the interior of a noncircular ellipse. For every
$a\in\Omega$ and every $d\ne0$, no map in $\Ead$ satisfies
\[
 \int_\Omega|\nabla u|=\Sigma(u).
\]
Thus no single-vortex class can witness property $(\mathrm P_1)$ on a
noncircular ellipse.
\end{corollary}

\section{Proof intuition}

The largest disk centered at the vortex is already enough to account for the
entire sharp lower bound. Every centered circle inside that disk sees degree
$d$, and hence at least $2\pi|d|$ of tangential variation. Integrating over
the radius gives $2\pi|d|r_a$, exactly the source value
\eqref{eq:source-value}.

If a minimizer exists, there is no energy left for either radial variation
inside the disk or any variation outside it. The minimizer is therefore an
angular loop inside the disk and locally constant outside. Its nonzero degree
can change only at directions where the inscribed circle actually touches the
domain boundary. The contact set must have positive length. A connected
real-analytic boundary which shares a positive-length set with a circle must
be that circle. The unresolved part of Open Problem 1 is whether an ellipse
could have an extremizer with a more complicated singularity distribution.

\section{Proof}

We use standard Sobolev slicing and trace facts in the form presented, for
example, in \cite{Leoni}.

\begin{lemma}[sharp polar bound and its equality consequences]\label{lem:polar}
Let $\Omega$ be bounded and simply connected, $a\in\Omega$, $d\ne0$, and
$r=r_a$. Every $u\in\Ead$ satisfies
\begin{equation}\label{eq:polar-bound}
 \int_\Omega|\nabla u|\ge 2\pi|d|r.
\end{equation}
If equality holds, then
\begin{align}
 \nabla u&=0 &&\text{a.e. on }\Omega\setminus B_r(a),
 \label{eq:outside-zero}\\
 u(a+\rho e^{i\theta})&=q(e^{i\theta})
 &&\text{for a.e. }(\rho,\theta)\in(0,r)\times(0,2\pi),
 \label{eq:angular-form}
\end{align}
for some $q\in W^{1,1}(\Sone;\Sone)$ of degree $d$.
\end{lemma}

\begin{proof}
For almost every $\rho\in(0,r)$, the slice
$u_\rho(\theta)=u(a+\rho e^{i\theta})$ belongs to
$W^{1,1}(\Sone;\Sone)$. The factor $e^{i\varphi}$ in the definition of
$\Ead$ has a real-valued phase on almost every slice and hence degree zero.
Thus $\deg u_\rho=d$. The length-degree inequality gives
\begin{equation}\label{eq:length-degree}
 \int_0^{2\pi}|\partial_\theta u(\rho,\theta)|\dd\theta
 \ge 2\pi|d|
 \quad\text{for a.e. }\rho\in(0,r).
\end{equation}
Using polar coordinates in $B_r(a)\subset\Omega$,
\begin{align*}
 \int_\Omega|\nabla u|
 &\ge \int_{B_r(a)}|\nabla u|\\
 &=\int_0^r\!\int_0^{2\pi}
 \sqrt{\rho^2|\partial_\rho u|^2+|\partial_\theta u|^2}
 \dd\theta\dd\rho\\
 &\ge\int_0^r\!\int_0^{2\pi}|\partial_\theta u|
 \dd\theta\dd\rho
 \ge 2\pi|d|r.
\end{align*}
If equality holds, the energy on $\Omega\setminus B_r(a)$ vanishes, proving
\eqref{eq:outside-zero}. Equality in the pointwise middle inequality forces
$\partial_\rho u=0$ almost everywhere in the disk. The ACL/slicing
representative is therefore independent of $\rho$, which gives
\eqref{eq:angular-form}; its degree is $d$ by the preceding slice argument.
\end{proof}

\begin{proof}[Proof of Theorem~\ref{thm:main}]
Suppose first that $u\in\Ead$ attains $m_{a,d}$. Equations
\eqref{eq:source-value} and \eqref{eq:polar-bound} put us in the equality
case of Lemma~\ref{lem:polar}.

Let
\[
 K=\{e^{i\theta}\in\Sone:a+r e^{i\theta}\in\partial\Omega\}
\]
be the set of contact directions between the centered inscribed circle and
the domain boundary. We claim that $q$ in \eqref{eq:angular-form} is locally
constant on $\Sone\setminus K$. Indeed, if $e^{i\theta_0}\notin K$, then
$x_0=a+r e^{i\theta_0}$ is an interior point of $\Omega$. A small
neighborhood of $x_0$ is contained in $\Omega$ and is cut into two pieces by
an arc of $\partial B_r(a)$. On its exterior piece, \eqref{eq:outside-zero}
makes $u$ locally constant. Uniqueness of the $W^{1,1}$ trace across the
internal smooth arc forces the trace $q$ on the interior piece to have the
same constant value on a neighborhood of $\theta_0$.

Consequently $q'=0$ almost everywhere on $\Sone\setminus K$. On the other
hand, the degree formula for the absolutely continuous circle map $q$ gives
\[
 2\pi d=\int_0^{2\pi} q\wedge q'\dd\theta.
\]
Since $d\ne0$, the derivative $q'$ cannot be supported on a null set. Hence
$K$ has positive one-dimensional measure.

Consider the real-analytic function
$F(x)=|x-a|^2-r^2$. Its restriction to the connected real-analytic
one-manifold $\partial\Omega$ vanishes on the positive-measure set
$\partial\Omega\cap\partial B_r(a)$. The real-analytic identity principle
\cite{KrantzParks} implies that $F$ vanishes identically on
$\partial\Omega$. Thus $\partial\Omega=\partial B_r(a)$; boundedness and
simple connectivity imply $\Omega=B_r(a)$.

Conversely, if $\Omega=B_r(a)$, direct integration gives
\[
 \int_\Omega|\nabla w_{a,d}|
 =\int_0^r\!\int_0^{2\pi}\frac{|d|}{\rho}\rho
 \dd\theta\dd\rho=2\pi|d|r=m_{a,d}.
\]
So $w_{a,d}$ attains the infimum.
\end{proof}

\begin{proof}[Proof of Corollary~\ref{cor:ellipse}]
The boundary of an ellipse is connected and real analytic. A noncircular
ellipse is not a disk centered at any interior point. Theorem~\ref{thm:main}
therefore excludes attainment in every $\Ead$.
\end{proof}

\section{Verification and limitations}

The proof has been audited at the following pressure points.
\begin{itemize}
 \item The singularity $1/|x-a|$ is integrable in two dimensions.
 \item A $W^{1,1}$ phase has degree zero on almost every centered slice, so
       the degree used in \eqref{eq:length-degree} is stable throughout
       $\Ead$.
 \item Equality in the integral chain forces separately zero exterior energy
       and zero radial derivative; no strict-convexity argument is used.
 \item The trace argument is local at an internal smooth arc and does not
       assume continuity of $u$ in two dimensions.
 \item Positive-measure contact, rather than a single tangency, is essential
       for the real-analytic identity step.
\end{itemize}

The result does \emph{not} prove that a noncircular ellipse fails
$(\mathrm P_1)$. A possible extremizer with multiple vortices or a general
singular Jacobian is not ruled out. This is why the packet is classified as a
substantial partial result.

\section{Bounded novelty check}

On 22 July 2026, bounded searches were made using arXiv/web indexes and the
local run registry. Queries combined arXiv:1606.01526, the paper title and
DOI, the three authors, the exact phrase ``non circular ellipse'', and the
terms ``single vortex'', ``attainment'', and properties $(\mathrm P_1)$ and
$(\mathrm P_3)$. The searches found the source article \cite{BMS}, its short
announcement, and unrelated uses of the label $(P_1)$, but no later paper
stating the theorem above or resolving the ellipse question. This is bounded
evidence only, not a guarantee of novelty.

\section{Human-review recommendation}

The theorem should receive expert review, with particular attention to the
local trace step at $\partial B_r(a)\cap\Omega$ and to the stated scope. If
confirmed, the next mathematical target is to extend the equality analysis to
finitely many vortices, which is the most direct route toward the full ellipse
conjecture.

\begin{thebibliography}{9}
\bibitem{BMS}
H.~Brezis, P.~Mironescu, and I.~Shafrir,
\emph{Distances between classes in $W^{1,1}(\Omega;\mathbb S^1)$},
Calc. Var. Partial Differential Equations 57 (2018), Article 14;
arXiv:1606.01526. DOI: 10.1007/s00526-017-1280-z.

\bibitem{Leoni}
G.~Leoni, \emph{A First Course in Sobolev Spaces}, 2nd ed.,
Graduate Studies in Mathematics 181, American Mathematical Society, 2017.

\bibitem{KrantzParks}
S.~G.~Krantz and H.~R.~Parks,
\emph{A Primer of Real Analytic Functions}, 2nd ed., Birkh\"auser, 2002.
\end{thebibliography}

\end{document}
